\section{Model Architectures}\label{sec:model-architectures}
In this project, a \gls{cnn} and an ensemble of \glspl{fb-mlp} with single-segment inputs were trained per patient to predict seizures.
This section discusses their exact specifications.

\subsection{Convolutional Neural Network}\label{subsec:cnn-architecture}
The architecture of the \gls{cnn} was adapted from \textcite{muller_coherent_2022}, which used the \textit{nv1x16} topology originally proposed in \textcite{eberlein_convolutional_2018}.
It uses the raw \gls{eeg} signals per \gls{eeg} channel from a segment as inputs.
Until the final stages, it treats \gls{eeg} channels separately, computing features per channel.
Only at the end are the per-channel features combined to output a final score.

Note that in the following, \gls{eeg} channels refer to electrode-derived input signals, whereas \gls{cnn} channels refer to convolutional feature maps (tensor depth).

\Cref{fig:cnn-architecture} shows the exact architecture.
The inputs are passed through a combination of convolutional and pooling layers.
There are two distinct types of convolutional blocks (\cref{fig:special-cnn-blocks}):
\begin{itemize}
    \item \textbf{conv1} uses a convolutional layer, followed by a pooling layer with a relatively large kernel size, reducing the height (number of \gls{eeg} samples) rapidly.
    \item \textbf{conv2} blocks further reduce the height by using valid padding for \textit{two} convolutional layers and pooling again.
\end{itemize}
All convolution layers use L1 and L2 regularization with a penalty of $10^{-9}$.

First, four conv1 blocks are used to reduce the tensor height while increasing the \gls{cnn} channels (encoding).
These are followed by three conv2 blocks decreasing the tensor height further while also decreasing the number of \gls{cnn} channels (decoding).
Ultimately, the sample dimension (height) is reduced to 1 with 32 features per \gls{eeg} channel.
Only then are the per-channel features combined (flattened).
These pass through two final dense layers to create the output logit, which is then bounded between 0 and 1 with the sigmoid function, to create the final output.

\begin{figure}[!htbp]
    \centering

    \begin{subfigure}[t]{\textwidth}
        \centering
        \includegraphics{conv_blocks}
        \caption{Special CNN blocks}
        \label{fig:special-cnn-blocks}
    \end{subfigure}

    \vspace{1ex}

    \begin{subfigure}[t]{\textwidth}
        \centering
        \includegraphics[width=\linewidth]{CNN_architecture}
        \caption{Full CNN architecture}
        \label{fig:full-cnn-architecture}
    \end{subfigure}

    \caption[CNN architecture]{\textbf{CNN architecture}.\\
    A segment's EEG data is processed from left to right. The volumes represent the input tensor and the outputs of layers. In between the volumes, the network's layer's properties are shown.\\
    \textit{k}: kernel size, \textit{f}: \#filters, \textit{p}: pool size.}
    \label{fig:cnn-architecture}
\end{figure}

Since the height is reduced further and further in deeper layers, the extracted features receive information from more samples, i.e., the receptive field of deeper neurons grows on the time scale.
Since the per-channel features are only combined after the final conv block, information from the \gls{eeg} channels is evaluated separately until the end.

To accommodate two differences in the data structure compared to \textcite{eberlein_convolutional_2018}, two adjustments were made to the architecture.
Here, 2 \gls{eeg} channels were used instead of 16.
This meant there were 8 times fewer nodes in the first flattened tensor.
For this reason, the number of nodes in the first dense layer was reduced by a factor of 8: from 64 to 8.
The flat tensors thus had a height of 64 → 8 → 1, instead of 512 → 64 → 1.

Secondly, because the sampling rate of $\sim$\SI{207}{\Hz} here differed slightly from the sampling rate of \SI{200}{\Hz}, the number of input samples from a segment differed too (3105 vs.\ 3000).
Thus, to arrive at a height of 1 at the final conv layer, it was necessary to reduce this dimension slightly more.
For this reason, the pool padding of the second conv2 layer was changed from same to valid.
This minor change was enough to fully accommodate the sampling frequency difference.

\subsection{Ensemble of Feature-Based Multilayer Perceptrons}\label{subsec:ensemble-architecture}
The second model was as an ensemble using \glspl{fb-mlp} as base learners, also adapted from \textcite{muller_coherent_2022}, where it had shown strong performance.
In contrast to architectures operating directly on the raw signal, this model relied on the hand-crafted features.
The input feature set comprised the spectral band powers and the correlation coefficient (\cref{subsec:feature-extraction}), which have previously been shown to possess forecasting utility in a cycle-based framework \parencite{yang_seizure_2024}.
This design enabled the model to incorporate both channel-wise feature characteristics and their inter-channel statistical dependencies.

Each individual \textbf{\gls{fb-mlp}} (\cref{fig:fb-mlp}) received z-score normalized features as input.
The network consisted of three fully connected hidden layers with 16, 8, and 4 nodes, respectively, each employing the \gls{relu} activation function.
Batch normalization was applied after the input layer and after each hidden layer to stabilize and accelerate training.
The output layer comprised a single unit with a sigmoid activation function to produce a probabilistic output.

\begin{figure}[p]
    \centering

    \begin{subfigure}[t]{\textwidth}
        \centering
        \includegraphics[width=\linewidth]{FB_MLP}
        \caption{Ensemble \textbf{input layer} and a single \textbf{\gls{fb-mlp}} with layers and their corresponding number of nodes.}
        \label{fig:fb-mlp}
    \end{subfigure}

    \vspace{1cm}

    \begin{subfigure}[t]{0.7\textwidth}
        \centering
        \includegraphics[width=\linewidth]{ensemble_architecture}
        \caption{Ensemble}
        \label{fig:ensemble-architecture}
    \end{subfigure}

    \vspace{0.4cm}
    
    \caption[Architecture of the ensemble of FB-MLPs]{Architecture of the ensemble of FB-MLPs.\\
    A segment's correlation coefficient and bandpowers are processed from left to right.}
    \label{fig:ensemble}
\end{figure}

The \textbf{ensemble} (\cref{fig:ensemble-architecture}) consisted of 50 \glspl{fb-mlp} that shared a common input layer but were initialized with distinct random weight configurations.
The final prediction was obtained by averaging the scalar outputs of the individual ensemble members.
While \textcite{muller_coherent_2022} employed an ensemble size of 100, their study considered
 a substantially more diverse set of input features, which allowed for greater heterogeneity among ensemble members.
Given the more constrained feature space in the present work, a smaller ensemble size of 50 base learners was selected as a trade-off to limit computational and memory requirements during training, while retaining the benefits of ensembling.


\subsection{Model Training} \label{subsec:model-training}
For model training, only \textbf{preictal} and \textbf{interictal} segments were included. 

\textbf{Inter-pre} and \textbf{intervention} segments were excluded because they occur immediately prior to seizure onset and may therefore share electro-physiological characteristics with preictal segments. 
If included, these segments—while still annotated with the label 0 (non-preictal)—could induce the model to learn misleading patterns.

Similarly, \textbf{postictal} and \textbf{inter-post} segments were omitted, as they may fall within the preictal window of seizures that did not meet the validity criteria and could consequently contain preictal-like features. 
Moreover, the postictal period is characterized by substantial alterations in brain state, leading to marked changes in signal dynamics that may confound model learning and degrade performance \parencite{muller_coherent_2022}.

As mentioned in \cref{subsec:dataset-partitioning}, there was a large class imbalance with scarcity of preictal segments.
To avoid wasting resources training on more interictal segments than necessary, the number of interictal segments was subsampled to a maximum of 20 times the number of preictal segments.
The class weights were adjusted to the final ratio of preictal vs.\ interictal segments, so that the models were not biased towards focusing mostly on learning to predict the majority (interictal) class.
For a given patient, all models received the same subsampled segments for training, so that results were more comparable.

As in \textcite{eberlein_convolutional_2018}, the \textbf{\glspl{cnn}} were trained for \textit{50 epochs} with \textit{ADAM optimization}, a \textit{learning rate of \num{0.001}}, and \textit{binary cross-entropy} as a loss function.

To train the \textbf{\glspl{fb-mlp}}, before subsampling, a \textit{z-score normalizer} was fit to the features of the entire training set.
This same scaling was then always applied when this model was later used, including for the test set.
The \glspl{fb-mlp} were then trained for \textit{200 epochs} using \textit{stochastic gradient descent} for optimization, a \textit{learning rate of \num{0.0001}}, and \textit{binary cross-entropy} as a loss function.
Even though \textcite{muller_coherent_2022} used 500 epochs for training, here, the performance could not be seen to improve after 200 epochs, which is likely also due to the smaller input dimension, requiring less fitting.
